\section{Introduction and Problem Understanding}

\subsection{System Identification}

This document specifies the design of \textbf{FalconVQA}, short for Fast Augmented
Language-based CONversational Video Question Answering. It is a retrieval and question
answering system for recorded video, developed with surveillance footage as its
primary application domain. The name records the four properties the system is
designed around: retrieval that is \textbf{fast} enough to be interactive, an index
\textbf{augmented} beyond the raw signal with structured descriptions produced by vision
and audio models, a \textbf{language-based} interface in which both the index and the
query are natural language, and a \textbf{conversational} access path in which an
operator or an autonomous agent asks questions of a recording and receives cited
answers rather than a list of files.

\subsection{Problem Statement}

A natural language front end that extracts the relevant footage from a given recorded
video by description alone.

\subsubsection*{Background}

Surveillance footage is generally recorded onto network video recorders and other
digital storage devices. Finding a particular piece of footage is carried out by
entering the date and time of the incident and reviewing what comes back, which is
slow. The recording exists as a continuous timeline indexed by nothing but a clock, so
the date and time are the only handles the operator has.

That procedure works only when the time of the incident is already known. In practice it
usually is not. What is available instead is a description of the event, obtained from a
witness or from a report: a man in a yellow shirt left a bag near the checkout. No stage
of the recording chain produces any account of what the footage contains, so a
description of this kind cannot be used as a query at all.

\subsubsection*{Challenge}

To develop a simple natural language front end in which the operator types a description
of the event into a prompt box, and the matching footage is extracted and played back.

Two things stand between the description and the playback. The first is that the footage
carries no description to match against, so one has to be derived. The second is that
general-purpose video understanding tools do not derive it, because they are built for
edited media and rely on three properties surveillance footage does not have:

\begin{itemize}
  \item \textbf{Shot boundaries.} A fixed camera produces none.
  \item \textbf{Speech.} A substantial proportion of surveillance footage is silent.
  \item \textbf{Visually salient events.} Events of interest are typically brief and
        visually unremarkable.
\end{itemize}

Applied to such footage, these tools either return the entire recording as a single
segment or subdivide it into a large number of near-identical segments. Neither result
is searchable.

The problem this project addresses is therefore twofold: derive a descriptive
representation of footage that has little intrinsic structure, and make the footage
retrievable through that representation. The operator states what is being looked for in
plain language and receives the corresponding moment, timestamped and ready to play.

\subsection{Project Objectives}

The system is required to deliver the following:

\begin{enumerate}[label=O-\arabic*.,leftmargin=3.0em]
  \item A processing pipeline that accepts a recording and renders it searchable by
        description.
  \item Semantic search returning ranked moments with timecodes, and question
        answering that returns a synthesised answer with references to the moments it
        was derived from.
  \item Video-level output that can be read in place of watching the footage: a
        summary, chapters, notable events, the people and objects that appear, and
        any anomalies.
  \item A multi-tenant web application in which uploading, searching, questioning and
        playback are performed in a single workspace, with each operator's recordings
        isolated from every other operator's.
  \item A documented and self-describing HTTP API, so that the same capabilities can
        be driven by an LLM agent rather than only by a human operator.
\end{enumerate}

\subsection{Document Organisation}

Section~2 presents the solution and the processing pipeline in outline. Section~3 states
the functional and non-functional requirements. Section~4 records the design
enhancements that distinguish the system from a conventional caption-and-embed pipeline.
Section~5 is the detailed design: each of the four pipeline stages, the models employed,
the parameters governing them and the outputs produced. Section~6 specifies the
architecture and technology stack, and Section~7 the database design across all four
stores. Section~8 describes the user interface. Section~9 states the acceptance criteria
and the test cases that establish them, followed by the plan under which those cases are
run: the environment, the test data, the entry and exit criteria, and what is automated.
Section~10 sets out the execution schedule and the distribution of work.
